\section*{Bab \ref{ch:b3:quantum-model-systems} --- Mekanika Kuantum untuk Kimiawan: Sistem Model}

\begin{solution}{exo:b3:quantum-model-systems:1}
$E_1 = h^2/8m_eL^2$. (a) $L = \qty{1.00}{nm}$: $E_1 = \qty{6.02e-20}{J} =
\qty{0.376}{eV}$; $\Delta E_{12} = 3E_1$ dan $\lambda = hc/3E_1 =
\qty{1099}{nm}$, di inframerah dekat. (b) $L = \qty{0.50}{nm}$: $E_1$ empat
kali lebih besar, \qty{2.41e-19}{J} $= \qty{1.50}{eV}$, dan $\lambda =
\qty{275}{nm}$, di ultraviolet.
\end{solution}

\begin{solution}{exo:b3:quantum-model-systems:2}
$n_x^2 + n_y^2 + n_z^2 = 3$ (1,1,1): $g = 1$; $6$ (2,1,1 dan permutasinya):
$g = 3$; $9$ (2,2,1): $g = 3$; $11$ (3,1,1): $g = 3$; $12$ (2,2,2): $g = 1$.
\end{solution}

\begin{solution}{exo:b3:quantum-model-systems:3}
$[\hat x^2,\hat p]f = -\iu\hbar x^2f' + \iu\hbar(x^2f)' = 2\iu\hbar xf$, sehingga
$[\hat x^2,\hat p] = 2\iu\hbar\hat x$. $[\hat p^2,\hat x]f = -\hbar^2[(xf)'' -
xf''] = -2\hbar^2f'$, sehingga $[\hat p^2,\hat x] = -2\iu\hbar\hat p$.
\end{solution}

\begin{solution}{exo:b3:quantum-model-systems:4}
$E_J/hc = BJ(J+1)$: 0, 21.19, 63.56 dan \qty{127.12}{cm^{-1}}, dengan $g = 1$,
3, 5, 7. $I = h/(8\pi^2cB) = \qty{2.643e-47}{kg.m^2}$ (dengan $c$ dalam
\unit{cm/s}).
\end{solution}

\begin{solution}{exo:b3:quantum-model-systems:5}
$|\psi_n|^2$ simetris terhadap $L/2$, sehingga $\langle x\rangle = L/2$. Dengan
$\sin^2u = \frac12(1 - \cos2u)$ dan dua kali integrasi parsial,
$\langle x^2\rangle = L^2\big(\frac13 - \frac{1}{2n^2\pi^2}\big)$. Untuk $n = 1$,
$\langle x^2\rangle = 0.283L^2$, lebih kecil daripada nilai klasik $L^2/3$:
keadaan dasar menumpuk di tengah. Ketika $n$ bertambah, nilai klasik itu
diperoleh kembali.
\end{solution}

\begin{solution}{exo:b3:quantum-model-systems:6}
$1\,\unit{cm^{-1}} = \qty{0.011963}{kJ/mol}$. ZPE(\ce{H2}) $= 2200.6 \times
0.011963 = \qty{26.33}{kJ/mol}$, ZPE(\ce{D2}) $= \qty{18.63}{kJ/mol}$; selisihnya
\qty{7.69}{kJ/mol}. \omterm{def:b3:quantum-model-systems:oscillator}{Tetapan gayanya} sama, sehingga $\tilde\omega \propto \mu^{-1/2}$
dan rasio yang diharapkan adalah $\sqrt{\mu_D/\mu_H} = 1.4137$ (massa atom
2.01410 dan 1.00783); rasio terukurnya 1.4127, dengan selisih 0.07\,\% (selisih
kecil itu berasal dari elektron, yang tidak mengikuti inti dengan sempurna).
\end{solution}

\begin{solution}{exo:b3:quantum-model-systems:7}
$\int_0^Lx^2(L - x)^2\,\dd x = L^5/30$, sehingga $N = \sqrt{30/L^5}$. Maka
\[
\langle\psi_1|\phi\rangle = \sqrt{\tfrac2L}\sqrt{\tfrac{30}{L^5}}\int_0^Lx(L -
x)\sin\tfrac{\pi x}{L}\,\dd x = \frac{\sqrt{60}}{L^3}\cdot\frac{4L^3}{\pi^3} =
\frac{4\sqrt{60}}{\pi^3} = 0.99928 .
\]
Parabola itu terdiri atas 99.86\,\% keadaan dasar (kuadrat tumpang-tindihnya).
\end{solution}

\begin{solution}{exo:b3:quantum-model-systems:8}
$\langle f|\hat pg\rangle = -\iu\hbar\int f^*g'\,\dd x = -\iu\hbar[f^*g] +
\iu\hbar\int f^{*\prime}g\,\dd x = \int(-\iu\hbar f')^*g\,\dd x = \langle\hat
pf|g\rangle$, karena suku dalam kurung siku lenyap di kedua ujung. Untuk
$\hat x\hat p$, dengan memakai bahwa $\hat x$ dan $\hat p$ masing-masing
\omterm{def:b3:quantum-model-systems:hermitian}{Hermitian}: $\langle f|\hat x\hat pg\rangle = \langle\hat xf|\hat pg\rangle = \langle\hat p\hat
xf|g\rangle$, dan $\hat p\hat x = \hat x\hat p - \iu\hbar \ne \hat x\hat p$: tidak
\omterm{def:b3:quantum-model-systems:hermitian}{Hermitian} (bentuk simetrisnya, $\frac12(\hat x\hat p + \hat p\hat x)$, \omterm{def:b3:quantum-model-systems:hermitian}{Hermitian}).
\end{solution}

\begin{solution}{exo:b3:quantum-model-systems:9}
$L = 6 \times 139.7 = \qty{838}{pm}$, $N = 6$: $\lambda = 8m_ecL^2/(7h) =
\qty{331}{nm}$. Serapannya berada di ultraviolet: heksatriena tidak berwarna.
\end{solution}

\begin{solution}{exo:b3:quantum-model-systems:10}
$\kappa = \sqrt{2m(V_0 - E)}/\hbar$ dengan $V_0 - E = \qty{0.20}{eV}$:
$\kappa_H = \qty{9.82e10}{m^{-1}}$, $\kappa_D = \qty{1.388e11}{m^{-1}}$.
Faktor-faktor di depan saling meniadakan dalam rasio: $T_H/T_D = \eu^{2a(\kappa_D - \kappa_H)} =
\eu^{4.06} = 58$, yaitu nilai eksaknya: penghalang itu tebal untuk keduanya
($\kappa a = 4.9$ dan 6.9).
\end{solution}

\begin{solution}{exo:b3:quantum-model-systems:11}
Tingkat energi $E_m = m^2\hbar^2/2m_eR^2$: $m = 0$ menampung dua elektron,
$m = \pm1$ empat. Tingkat terisi tertinggi adalah $|m| = 1$ dan tingkat kosong
terendah $|m| = 2$: $\Delta E =
3\hbar^2/2m_eR^2 = \qty{9.38e-19}{J}$,
$\lambda = \qty{212}{nm}$, di ultraviolet, tempat serapan kuat benzena juga
berada.
\end{solution}

\begin{solution}{exo:b3:quantum-model-systems:12}
$\langle r\rangle = \frac{4\pi}{\pi a^3}\int_0^\infty r^3\eu^{-2r/a}\dd r =
\frac{4}{a^3}\cdot\frac{3!}{(2/a)^4} = \frac32a$, \qty{79.4}{pm} untuk $a = a_0$.
$\langle 1/r\rangle = \frac{4}{a^3}\cdot\frac{1!}{(2/a)^2} = \frac1a$, sehingga
$\langle V\rangle = -e^2/(4\pi\varepsilon_0a)$. Karena $E_1 =
-e^2/(8\pi\varepsilon_0a)$ (dari $a = 4\pi\varepsilon_0\hbar^2/\mu e^2$),
$\langle V\rangle = 2E_1$, dan energi kinetik rata-ratanya $-E_1$.
\end{solution}

\begin{solution}{pb:b3:quantum-model-systems:1}
\textbf{1.} Setiap atom dari $j$ atom itu menyumbang satu orbital p; $j - 1$
atom karbon dan satu nitrogen masing-masing menyumbang satu elektron, dan
nitrogen yang lain menyumbang pasangan elektron bebasnya, karena muatan kation
tersebar di sepanjang rantai: $N = j + 1$, yaitu 6, 8 dan 10.
\textbf{2.} Enam elektron mengisi $n = 1$, 2, 3: HOMO $n = 3$, LUMO $n = 4$.
\textbf{3.} $\Delta E = E_{N/2+1} - E_{N/2} = (N+1)h^2/8mL^2$ dan $\lambda =
hc/\Delta E = 8mcL^2/h(N+1)$.
\textbf{4.} $L_0 = 4l$, $6l$, $8l$: 559, 838 dan \qty{1118}{pm}.
\textbf{5.} 147, 257 dan \qty{374}{nm}.
\textbf{6.} Semuanya jauh terlalu pendek (dengan faktor 3.6, 2.3 dan 1.9):
kotak dalam model itu terlalu kecil, karena $\lambda \propto L^2$; elektron
bergerak sepanjang jarak yang lebih jauh daripada rantai di antara kedua
nitrogen.
\textbf{7.} Dengan titik asal di pusat, $\psi_n \propto \cos(n\pi u/L)$ untuk
$n$ ganjil dan $\sin(n\pi u/L)$ untuk $n$ genap ($u = x - L/2$): fungsi genap
dan fungsi ganjil dari $u$.
\textbf{8.} Jika $n + n'$ genap, $\psi_n$ dan $\psi_{n'}$ mempunyai paritas
yang sama; hasil kali keduanya dengan $u$ ganjil dan integralnya pada selang
yang simetris bernilai nol.
\textbf{9.} HOMO $N/2$ dan LUMO $N/2 + 1$ memberikan $n + n' = N + 1$, ganjil:
selalu diizinkan.
\textbf{10.} Zat warna kedua: HOMO 4, LUMO 5. $4 \to 6$: jumlahnya genap,
terlarang. $3 \to 5$: jumlahnya genap, terlarang.
\textbf{11.} Transisi diizinkan berikutnya dari $n = 4$ adalah $4 \to 7$,
dengan $\Delta E$ lebih besar $(49 - 16)/(25 - 16) = 3.67$ kali: pada
$\lambda/3.67$, yaitu \qty{164}{nm} untuk zat warna kedua, jauh di ultraviolet.
\textbf{12.} Aturan itu hanya bergantung pada simetri potensial terhadap
pusat rantai, yang juga dimiliki zat warna simetris yang sebenarnya.
\textbf{13.} $E = hc/\lambda = \qty{3.294e-19}{J} = \qty{2.06}{eV}$.
\textbf{14.} $hc\tilde\omega_e = \qty{0.371}{eV}$.
\textbf{15.} $2hcB = \qty{4.21e-22}{J} = \qty{2.63}{meV}$.
\textbf{16.} $kT = \qty{25.7}{meV}$.
\textbf{17.} Hanya tingkat-tingkat rotasi ($2hcB \ll kT$); eksitasi vibrasi
($15kT$) dan elektronik ($80kT$) praktis tidak terpopulasi.
\textbf{18.} Elektronik: tampak dan ultraviolet; vibrasi: inframerah;
rotasi: gelombang mikro (dan inframerah jauh).
\textbf{19.} $L = \sqrt{\lambda h(N+1)/8mc}$ dengan $N = 8$: $L = \qty{1283}{pm}$.
\textbf{20.} $\delta = (1283 - 838)/2 = \qty{222}{pm}$.
\textbf{21.} $L = 559 + 445 = \qty{1003}{pm}$, $\lambda = \qty{474}{nm}$, 9.5\,\%
di bawah 524 nm.
\textbf{22.} $L = 1118 + 445 = \qty{1562}{pm}$, $\lambda = \qty{732}{nm}$, 2.9\,\%
di atas 711 nm.
\textbf{23.} Ke dalam cincin-cincin aromatik yang membawa atom nitrogen:
sistem $\pi$ tidak berakhir di atom nitrogen, dan kotak efektifnya mencakup
sebagian dari setiap cincin.
\textbf{24.} $\delta/l = 222/139.7 = 1.6$: \textbf{kotak itu memanjang sekitar
\qty{222}{pm}, satu setengah ikatan, melewati setiap nitrogen}, dan dengan
satu panjang itu saja model ini mereproduksi ketiga warna dengan selisih
kurang dari 10\,\%.
\end{solution}
